\section*{Appendix A: The Beachhead Definition Canvas}

Complete this canvas for exactly 1 proposed first market. Replace adjectives with observable evidence wherever possible:

\begin{table}[htbp]
\centering
\small
\begin{tabularx}{\columnwidth}{l X}
\toprule
\textbf{Question} & \textbf{Answer / Evidence Status} \\
\midrule
Market context & Category and overall demand system \\
Target user & Professional role and day-to-day context \\
Economic buyer & Who authorizes payment and budget \\
Triggering event & What changed that makes action urgent \\
Important job & Exact progress or failure avoided \\
Cost of delay & Quantifiable financial or career penalty \\
Default alternative & Current workaround or competing tool \\
Why alternative fails & Root cause of dissatisfaction \\
Prerequisites & Starting capability and technical baseline \\
Initial complete offer & Exact scope of initial deliverable \\
Reason to try us now & Wedge mechanism and adoption incentive \\
Where buyer searches & Communities, networks, or trusted peers \\
Existing proof & Work samples, rubrics, or references \\
Ethical limits & Exclusions, confidentiality, regulations \\
Price hypothesis & Realized fee and payment structure \\
Cost hypothesis & Direct variable cost and founder hours \\
First experiment & Scope, duration, and participant cap \\
Decision thresholds & Criteria to repeat, redesign, or stop \\
Adjacent groups & Next segments and transferable assets \\
\bottomrule
\end{tabularx}
\caption{The 1-page beachhead definition canvas.}
\label{tab:ch7-canvas}
\end{table}

\textbf{One-sentence thesis template:} ``For [identified buyer] facing [real trigger], existing choices fail at [specific job] because [observed mechanism]. We can reach them through [real channel] and deliver [small complete outcome] with [credible proof], at economics that [test]. Success would produce [transferable asset].''

\section*{Appendix B: Segment Discovery and Anti-Segment Map}

For each candidate market group, record observable behavioral differences:

\begin{table*}[t]
\centering
\small
\begin{tabularx}{\textwidth}{l p{2.4cm} p{1.8cm} p{2.4cm} p{2.0cm} p{2.4cm} X}
\toprule
\textbf{Candidate} & \textbf{Trigger and Job} & \textbf{Payer} & \textbf{Costly Workaround} & \textbf{Reachability} & \textbf{Why Underserved} & \textbf{Anti-Customer} \\
\midrule
Segment A & Deployment deadline; reliability audit & Self-pay engineer & Senior code reviews; manual test runs & Specialized tech communities & Incumbents lack deep technical feedback & Complete beginners seeking placement \\
Segment B & Semester elective; applied lab need & College department & Guest lectures; ad hoc faculty labs & Academic leadership network & Rigid syllabi and non-applied curriculum & Large colleges with no autonomous credits \\
Segment C & Enterprise feature launch; QA failure & Engineering VP & Generic corporate training courses & Targeted executive introductions & Courses lack custom production context & Teams requiring custom on-premise IP \\
\bottomrule
\end{tabularx}
\caption{Segment discovery and anti-segment disqualification map.}
\label{tab:ch7-anti-segment-map}
\end{table*}

Do not treat the existence of a complaint as proof of buyer willingness to pay. Require observed sacrifice or purchasing evidence before substantial investment.

\section*{Appendix C: The Hard-Gate and Evidence-Adjusted Comparison}

First apply \textbf{pass / redesign / defer} to the 6 non-negotiable feasibility gates from Section 38. For survivors only, use a 0 to 3 descriptive scale: 0 contradicted; 1 plausible but unsupported; 2 supported by direct preliminary evidence; 3 supported by repeated relevant evidence. Record evidence confidence separately.

\begin{table*}[t]
\centering
\small
\begin{tabularx}{\textwidth}{l p{2.6cm} p{2.6cm} p{2.6cm} X}
\toprule
\textbf{Criterion} & \textbf{Candidate A (D2C)} & \textbf{Candidate B (Colleges)} & \textbf{Candidate C (Enterprise)} & \textbf{Strongest Evidence / Missing Proof} \\
\midrule
Acute, consequential job & Score: 2 (observed) & Score: 1 (inferred) & Score: 2 (observed) & Verified deployment deadlines vs.\ general syllabus upgrades \\
Buyer authority and budget & Score: 2 (self-pay) & Score: 1 (committee) & Score: 2 (manager) & Personal card purchase vs.\ formal semester procurement \\
Repeated access to buyers & Score: 2 (communities) & Score: 1 (outreach) & Score: 1 (introductions) & Practitioner meetups vs.\ institutional sales calls \\
Meaningful difference & Score: 3 (deep review) & Score: 2 (rigor) & Score: 2 (custom task) & Rigorous technical critique vs.\ mass video lectures \\
Delivery predictability & Score: 2 (remote studio) & Score: 1 (travel/labs) & Score: 1 (bespoke IP) & Controlled cohort studio vs.\ college administrative load \\
Economic viability & Score: 2 (positive cash) & Score: 1 (slow payment) & Score: 3 (high fee) & Upfront tuition vs.\ net-90 institutional invoicing \\
Fast learning in runway & Score: 3 (3-week cycle) & Score: 1 (semester) & Score: 1 (quarterly) & Rapid feedback loop vs.\ long academic approval cycles \\
Bounded downside risk & Score: 3 (clear terms) & Score: 2 (contracts) & Score: 1 (security NDAs) & Standard refund terms vs.\ enterprise liability clauses \\
Portable proof & Score: 3 (code artifacts) & Score: 2 (grades/labs) & Score: 1 (confidential) & Public technical defense vs.\ proprietary client systems \\
Reusable asset transfer & Score: 2 (curriculum) & Score: 3 (course packs) & Score: 2 (case studies) & Modular evaluation harnesses transfer across sectors \\
\bottomrule
\end{tabularx}
\caption{Evidence-adjusted scorecard comparing 3 candidate entry positions.}
\label{tab:ch7-scorecard}
\end{table*}

\textbf{Interpretation rule:} Do not automatically choose the highest total. Inspect fatal weaknesses, confidence levels, founder constraints, and the single experiment that could resolve the most consequential disagreement.

\section*{Appendix D: A 20-Question Beachhead Interview Guide}

Ask about a recent real buying or problem-solving event. Do not read this mechanically as a survey, and obtain permission before recording:

\subsection*{Context and Trigger}
\begin{enumerate}
    \item What specific work were you responsible for when this issue arose?
    \item What changed internally or externally to make this urgent right then?
    \item Who else in your organization or team was affected by the outcome?
    \item What external deadline or milestone forced a prompt decision?
\end{enumerate}

\subsection*{Current Solution and Workarounds}
\begin{enumerate}
    \setcounter{enumi}{4}
    \item What did you try first when you realized you faced this problem?
    \item Which tools, suppliers, colleagues, or manual workarounds did you deploy?
    \item What did that temporary workaround cost in time, money, or stress?
    \item Which specific tasks or edge cases did existing vendors completely fail to resolve?
\end{enumerate}

\subsection*{Decision and Alternatives}
\begin{enumerate}
    \setcounter{enumi}{8}
    \item Which alternative products or courses did you evaluate before deciding?
    \item Why did you reject the options you ultimately did not select?
    \item Who had the actual authority to approve spending and sign off on the solution?
    \item Was dedicated budget already available, or did funds have to be reallocated?
    \item What non-financial costs (such as reputation or disruption) mattered most?
    \item What would have happened if you had simply postponed the decision by 3 months?
\end{enumerate}

\subsection*{Proof and Access}
\begin{enumerate}
    \setcounter{enumi}{14}
    \item Where did you look for advice or trusted recommendations?
    \item Which kind of reference or credential would have made an unknown vendor credible?
    \item What operational risks or liabilities would have immediately ruled out a provider?
    \item What conditions made a low-risk paid pilot acceptable to your approver?
\end{enumerate}

\subsection*{Outcome and Next Job}
\begin{enumerate}
    \setcounter{enumi}{18}
    \item What was the measurable outcome after you implemented your choice?
    \item Have you faced a similar decision since, or advised a colleague in that exact position?
\end{enumerate}

After each interview, write separately: observed events, respondent reports, founder interpretations, and untested assumptions. Group responses by \emph{decision situation}, not merely by polite praise for your proposed offer.
